\documentclass[11pt,a4paper]{article}
% Preâmbulo comum. Não duplicar pacotes nos módulos.
\usepackage[utf8]{inputenc}
\usepackage[T1]{fontenc}
\usepackage[brazil]{babel}
\usepackage{amsmath,amssymb,amsthm,mathtools}
\usepackage{geometry}
\geometry{a4paper,margin=2.3cm}
\usepackage{microtype}
\usepackage{enumitem}
\usepackage{graphicx}
\usepackage{xcolor}
\usepackage{booktabs}
\usepackage{hyperref}
\hypersetup{colorlinks=true,linkcolor=blue!40!black,urlcolor=blue!50!black}

\setcounter{tocdepth}{2}
\numberwithin{equation}{section}

\theoremstyle{plain}
\newtheorem{teorema}{Teorema}[section]
\newtheorem{proposicao}[teorema]{Proposição}
\newtheorem{lema}[teorema]{Lema}
\newtheorem{corolario}[teorema]{Corolário}

\theoremstyle{definition}
\newtheorem{definicao}[teorema]{Definição}
\newtheorem{exemplo}[teorema]{Exemplo}
\newtheorem{exercicio}{Exercício}[section]

\theoremstyle{remark}
\newtheorem{observacao}[teorema]{Observação}
\newtheorem{contraexemplo}[teorema]{Contraexemplo}
\newtheorem{intuicao}[teorema]{Intuição}

\newcommand{\R}{\mathbb{R}}
\newcommand{\N}{\mathbb{N}}
\newcommand{\Z}{\mathbb{Z}}
\newcommand{\Q}{\mathbb{Q}}
\DeclareRobustCommand{\eps}{\varepsilon}
\DeclareRobustCommand{\vet}[1]{\vec{#1}}
\DeclareRobustCommand{\dif}{\mathrm{d}}
\DeclareRobustCommand{\abs}[1]{\left\lvert #1 \right\rvert}
\DeclareRobustCommand{\norm}[1]{\left\lVert #1 \right\rVert}
\DeclareMathOperator{\grad}{grad}
\DeclareMathOperator{\Div}{div}
\DeclareMathOperator{\rot}{rot}

\newcommand{\base}[1]{\par\medskip\noindent\textbf{Base para relembrar.} #1\par\medskip}
\newcommand{\cuidado}[1]{\par\medskip\noindent\textbf{Quando a hipótese falha.} #1\par\medskip}

\newcommand{\figura}[2]{%
  \begin{figure}[htbp]
  \centering
  \IfFileExists{#1}{%
    \includegraphics[width=0.84\textwidth]{#1}%
  }{%
    \fbox{\parbox{0.8\textwidth}{\centering Figura ainda não gerada:\\ \texttt{#1}}}%
  }
  \caption{#2}
  \end{figure}%
}

\newenvironment{esqueleto}{%
  \par\medskip\noindent\textbf{Esqueleto da demonstração.}\par\nopagebreak
  \begin{enumerate}[leftmargin=1.4em]
}{\end{enumerate}\medskip}

\title{Cálculo III: integrais de linha e campos conservativos}
\author{}
\date{}
\begin{document}
\maketitle
\tableofcontents

\section{Campos escalares e vetoriais}

Um campo associa a cada ponto de uma região uma grandeza. A grandeza pode ser um número ou um vetor. A integral de linha vai somar essa grandeza ao longo de uma curva: no caso escalar, pesando pelo comprimento; no caso vetorial, pesando pela componente tangente. Trabalho, massa de um fio e circulação são os três nomes físicos dessa soma.

\begin{definicao}
Seja $U\subset\R^n$ aberto. Um \emph{campo escalar} em $U$ é uma função $f\colon U\to\R$. Um \emph{campo vetorial} em $U$ é uma função $\vet{F}\colon U\to\R^n$. Escrevemos as componentes
\[
\vet{F}=(F_1,\ldots,F_n)
\]
e, no plano e no espaço, também $\vet{F}=(P,Q)$ e $\vet{F}=(P,Q,R)$.
\end{definicao}

O aberto garante vizinhança em cada ponto, que é o que as derivadas parciais e a regra da cadeia vão pedir.

\begin{exemplo}
Três campos com unidade mental.
\begin{enumerate}
\item A temperatura $T(x,y,z)$, em kelvin, é um campo escalar. O valor não tem direção.
\item A velocidade $\vet{v}(x,y,z)$ de um fluido, em metro por segundo, é um campo vetorial. A seta em cada ponto é a velocidade da partícula que passa por ali.
\item Se $T$ é de classe $C^1$, o gradiente $\nabla T$ é um campo vetorial, em kelvin por metro. Ele aponta para o crescimento mais rápido de $T$. O fluxo de calor, no modelo mais simples, aponta para o lado oposto: a direção de $-\nabla T$.
\end{enumerate}
Uma força $\vet{F}$, em newton, é um campo vetorial do mesmo tipo que a velocidade: o que muda é a unidade, não a geometria da integral.
\end{exemplo}

\begingroup\catcode`\_=12
\figura{../figuras/campo_vetorial.png}{À esquerda um campo conservativo; à direita $\vet{F}=(-y,x)$.}
\endgroup

À esquerda, as setas variam como um gradiente: não há giro líquido em volta de um ponto. À direita, $\vet{F}=(-y,x)$ é perpendicular ao raio vetor, porque
\[
\vet{F}(x,y)\cdot(x,y)=-yx+xy=0,
\]
e $\norm{\vet{F}}=\sqrt{x^2+y^2}$ é a distância à origem. O campo é tangente às circunferências e gira no sentido anti-horário. A circulação desse giro, na circunferência unitária, será $2\pi$.

\begin{intuicao}
O campo escalar responde à pergunta ``quanto tem aqui?''. O campo vetorial responde ``para onde e com que intensidade?''. A integral de linha do escalar soma ``quanto'' vezes ``quanto de fio''. A integral de linha do vetor soma só o que está alinhado com a curva.
\end{intuicao}

\section{Integral de linha de um escalar}

\begin{definicao}
Seja $f$ contínua num aberto que contém a imagem da curva $\vet{r}\colon[a,b]\to\R^n$, de classe $C^1$. A \emph{integral de linha escalar} de $f$ ao longo de $C=\vet{r}([a,b])$ é
\[
\int_C f\,\dif s
=\int_a^b f\bigl(\vet{r}(t)\bigr)\,\norm{\vet{r}'(t)}\,\dif t.
\]
Se $\vet{r}$ é apenas $C^1$ por pedaços, a integral é a soma das integrais nos pedaços.
\end{definicao}

O fator $\norm{\vet{r}'(t)}$ é a velocidade escalar. O produto $\norm{\vet{r}'(t)}\,\dif t$ é o elemento de comprimento: a integral não vê o nome do parâmetro, vê o comprimento. Em particular, $f=1$ recupera o comprimento de arco.

\begin{proposicao}
A integral de linha escalar não depende da orientação. Se $C^{-}$ é a mesma curva percorrida no sentido oposto e $f$ é contínua, então $\int_{C^{-}} f\,\dif s=\int_C f\,\dif s$.
\end{proposicao}

\begin{proof}
Seja $\vet{r}\colon[a,b]\to\R^n$ uma parametrização $C^1$ de $C$. A curva oposta admite
\[
\vet{r}^{-}(t)=\vet{r}(a+b-t),\qquad t\in[a,b].
\]
Então $(\vet{r}^{-})'(t)=-\vet{r}'(a+b-t)$ e $\norm{(\vet{r}^{-})'(t)}=\norm{\vet{r}'(a+b-t)}$. Logo
\[
\int_{C^{-}} f\,\dif s
=\int_a^b f\bigl(\vet{r}(a+b-t)\bigr)\,\norm{\vet{r}'(a+b-t)}\,\dif t.
\]
A substituição $u=a+b-t$ troca os limites e produz $\dif t=-\dif u$, portanto
\[
\int_b^a f\bigl(\vet{r}(u)\bigr)\,\norm{\vet{r}'(u)}\,(-\dif u)
=\int_a^b f\bigl(\vet{r}(u)\bigr)\,\norm{\vet{r}'(u)}\,\dif u.
\]
O sinal da derivada foi absorvido pela norma. A curva $C^1$ por pedaços reduz-se à soma dos pedaços.
\end{proof}

\begin{exemplo}[Massa de um fio]
A densidade linear $\delta$, em massa por comprimento, é um campo escalar sobre o fio. A massa é $\int_C \delta\,\dif s$: cada pedaço de comprimento $\dif s$ contribui com $\delta\,\dif s$, e o sentido em que se percorre o fio não muda a massa.

Tome $\delta(x,y)=2x$ ao longo de $\vet{r}(t)=(t,t^2)$, $t\in[0,1]$. Então $\vet{r}'(t)=(1,2t)$ e $\norm{\vet{r}'(t)}=\sqrt{1+4t^2}$. A massa é
\[
\int_0^1 2t\sqrt{1+4t^2}\,\dif t.
\]
Com $u=1+4t^2$, tem-se $\dif u=8t\,\dif t$, de modo que $2t\,\dif t=\dif u/4$. Os limites passam de $t=0$ a $t=1$ para $u=1$ a $u=5$:
\[
\int_1^5 \frac{1}{4}u^{1/2}\,\dif u
=\frac{1}{4}\cdot\frac{2}{3}\Bigl[u^{3/2}\Bigr]_1^5
=\frac{1}{6}\bigl(5^{3/2}-1\bigr)
=\frac{5\sqrt{5}-1}{6}.
\]
Se $x$ está em metros e $\delta$ em quilograma por metro, esse número é a massa em quilogramas.
\end{exemplo}

\section{Integral de linha de um campo vetorial}

\begin{definicao}
Seja $\vet{F}$ contínuo num aberto que contém a imagem de $\vet{r}\colon[a,b]\to\R^n$, de classe $C^1$. A \emph{integral de linha} de $\vet{F}$ ao longo de $C$ é
\[
\int_C \vet{F}\cdot\dif\vet{r}
=\int_a^b \vet{F}\bigl(\vet{r}(t)\bigr)\cdot\vet{r}'(t)\,\dif t.
\]
No plano, com $\vet{F}=(P,Q)$ e $\vet{r}=(x,y)$, isso é $\int_C P\,\dif x+Q\,\dif y$. No espaço, soma-se $R\,\dif z$.
\end{definicao}

Onde $\vet{r}'(t)\neq\vet{0}$, o vetor unitário tangente é $\vet{T}=\vet{r}'/\norm{\vet{r}'}$, e
\[
\vet{F}\cdot\vet{r}'\,\dif t
=(\vet{F}\cdot\vet{T})\,\norm{\vet{r}'}\,\dif t
=(\vet{F}\cdot\vet{T})\,\dif s.
\]
A integral vetorial só enxerga a componente tangencial. A componente normal à curva não trabalha e não circula.

\begin{proposicao}
A troca de orientação troca o sinal: $\int_{C^{-}}\vet{F}\cdot\dif\vet{r}=-\int_C\vet{F}\cdot\dif\vet{r}$.
\end{proposicao}

\begin{proof}
Com $\vet{r}^{-}(t)=\vet{r}(a+b-t)$,
\[
\vet{F}\bigl(\vet{r}^{-}(t)\bigr)\cdot(\vet{r}^{-})'(t)
=-\,\vet{F}\bigl(\vet{r}(a+b-t)\bigr)\cdot\vet{r}'(a+b-t).
\]
A mesma substituição $u=a+b-t$ carrega esse sinal negativo até o fim. Aqui não há norma para absorvê-lo: $\vet{T}$ vira $-\vet{T}$.
\end{proof}

\begin{intuicao}
Se $\vet{F}$ é uma força, $\int_C\vet{F}\cdot\dif\vet{r}$ é o trabalho dessa força ao longo do caminho, em joule quando a força está em newton e o caminho em metro. Ir e voltar pelo mesmo traço troca o sinal do trabalho. Empurrar a favor da força dá trabalho positivo; ser arrastado contra ela dá trabalho negativo.
\end{intuicao}

\begin{exemplo}[Circunferência unitária]
Seja $\vet{F}=(-y,x)$ e seja $C$ a circunferência $x^2+y^2=1$ orientada no sentido anti-horário,
\[
\vet{r}(t)=(\cos t,\sin t),\qquad t\in[0,2\pi].
\]
Então $\vet{r}'(t)=(-\sin t,\cos t)$ e $\vet{F}(\vet{r}(t))=(-\sin t,\cos t)$. O produto escalar é
\[
\vet{F}(\vet{r}(t))\cdot\vet{r}'(t)
=\sin^2 t+\cos^2 t=1.
\]
Portanto
\[
\int_C \vet{F}\cdot\dif\vet{r}
=\int_0^{2\pi} 1\,\dif t=2\pi.
\]
Geometricamente, $\vet{F}$ coincide com $\vet{r}'$ ao longo de $C$: a componente tangencial vale $1$ e a integral é o comprimento da circunferência. O sentido horário produz $-2\pi$.
\end{exemplo}

Esse valor já mostra que $\vet{F}=(-y,x)$ não pode ser um gradiente em $\R^2$. A conta geral está na seção seguinte: a integral de um gradiente numa curva fechada é zero, e $2\pi\neq 0$. O rotacional confirma o mesmo obstáculo sem integrar. Com $\vet{F}=(P,Q)=(-y,x)$,
\[
\partial_x Q-\partial_y P=\partial_x(x)-\partial_y(-y)=1-(-1)=2.
\]
O giro visto na figura é constante e igual a $2$, não uma ilusão de um caminho mal escolhido.

\section{Campo conservativo}

\base{Regra da cadeia: se $\phi$ é de classe $C^1$ e $\vet{r}=(x,y,z)$ é $C^1$, então $\frac{\dif}{\dif t}\phi(\vet{r}(t))=\nabla\phi(\vet{r}(t))\cdot\vet{r}'(t)$. Teorema fundamental do cálculo: se $g$ é $C^1$ em $[a,b]$, então $\int_a^b g'(t)\,\dif t=g(b)-g(a)$.}

\begin{definicao}
Seja $U\subset\R^n$ aberto e conexo, e seja $\vet{F}\colon U\to\R^n$ contínuo. Dizemos que $\vet{F}$ é \emph{conservativo} em $U$ quando existe $\phi\colon U\to\R$ de classe $C^1$ com $\vet{F}=\nabla\phi$ em $U$. A função $\phi$ é um \emph{potencial escalar} de $\vet{F}$.
\end{definicao}

A conexidade entra para que o potencial fale de qualquer par de pontos do domínio. A conta do teorema seguinte, numa curva já dada, usa só que $U$ é aberto.

\begin{teorema}
\label{thm:gradiente}
Seja $U\subset\R^n$ aberto e seja $\phi\colon U\to\R$ de classe $C^1$. Se $C$ é uma curva $C^1$ por pedaços em $U$, com extremidade inicial $A$ e extremidade final $B$, então
\[
\int_C \nabla\phi\cdot\dif\vet{r}=\phi(B)-\phi(A).
\]
Em particular, a integral depende de $A$ e de $B$, e não do caminho. Se $C$ é fechada, a integral é zero.
\end{teorema}

\begin{proof}
Seja primeiro $\vet{r}\colon[a,b]\to U$ de classe $C^1$, com $\vet{r}(a)=A$ e $\vet{r}(b)=B$. A composta $g(t)=\phi(\vet{r}(t))$ é de classe $C^1$ em $[a,b]$, porque $U$ é aberto e a imagem de $\vet{r}$, sendo compacta, fica a uma distância positiva do complementar de $U$. A regra da cadeia dá
\begin{align*}
g'(t)
&=\partial_x\phi(\vet{r}(t))\,x'(t)
+\partial_y\phi(\vet{r}(t))\,y'(t)
+\partial_z\phi(\vet{r}(t))\,z'(t)\\
&=\nabla\phi\bigl(\vet{r}(t)\bigr)\cdot\vet{r}'(t).
\end{align*}
No plano, o termo em $z$ não aparece. Pelo teorema fundamental do cálculo,
\[
\int_a^b \nabla\phi\bigl(\vet{r}(t)\bigr)\cdot\vet{r}'(t)\,\dif t
=g(b)-g(a)=\phi(B)-\phi(A).
\]
Se a curva é $C^1$ por pedaços, soma-se o resultado em cada pedaço. Os valores de $\phi$ nos pontos de emenda se cancelam, e sobram $\phi(B)$ e $-\phi(A)$. Numa curva fechada, $A=B$ e a diferença é zero.
\end{proof}

\begin{exemplo}
O campo $\vet{F}=(2xy,x^2)$ é o gradiente de $\phi(x,y)=x^2 y$, pois $\partial_x\phi=2xy$ e $\partial_y\phi=x^2$. Entre $A=(0,0)$ e $B=(1,3)$, qualquer caminho em $\R^2$ dá
\[
\int_C\vet{F}\cdot\dif\vet{r}=\phi(1,3)-\phi(0,0)=3.
\]
Conferência no segmento $\vet{r}(t)=(t,3t)$, $t\in[0,1]$: $\vet{F}(\vet{r}(t))=(6t^2,t^2)$, $\vet{r}'(t)=(1,3)$ e o produto escalar vale $9t^2$. Então $\int_0^1 9t^2\,\dif t=3$.
\end{exemplo}

\begin{exemplo}[Peso]
Perto do solo, o peso de uma massa $m>0$ é $\vet{F}=(0,0,-mg)$, com $g>0$ a aceleração da gravidade. Esse campo é o gradiente de $\phi(x,y,z)=-mgz$. O trabalho do peso entre as alturas $z_A$ e $z_B$ é
\[
\phi(B)-\phi(A)=-mg(z_B-z_A).
\]
Na descida, $z_B<z_A$ e o trabalho do peso é positivo. A energia potencial usual $V=mgz$ satisfaz $\vet{F}=-\nabla V$, e o mesmo trabalho se escreve $V(A)-V(B)$. As duas funções diferem pelo sinal: $\phi$ é o potencial do teorema, $V$ é a energia potencial da física.
\end{exemplo}

\section{A recíproca}

A ida está feita: gradiente implica integral dependente só dos extremos. A volta constrói o potencial a partir das integrais, desde que essas integrais não dependam do caminho.

\base{Num aberto de $\R^n$, cada componente conexa é conexa por caminhos, e dois pontos dela podem ser ligados por uma poligonal contida no aberto. Se as derivadas parciais de $\phi$ existem numa vizinhança de um ponto e são contínuas aí, então $\phi$ é diferenciável nesse ponto e o gradiente é o vetor das parciais.}

\begin{teorema}
\label{thm:reciproca}
Seja $U\subset\R^n$ aberto e conexo, e seja $\vet{F}=(P_1,\ldots,P_n)\colon U\to\R^n$ contínuo. Suponha que, para qualquer par de pontos de $U$ e quaisquer duas curvas $C^1$ por pedaços em $U$ com esses extremos, as integrais de linha de $\vet{F}$ coincidem. Então $\vet{F}$ é conservativo em $U$.
\end{teorema}

\begin{proof}
Fixe $A\in U$. Para cada ponto $X\in U$, escolha uma curva $C^1$ por pedaços $\gamma_X$ em $U$ ligando $A$ a $X$ e defina
\[
\phi(X)=\int_{\gamma_X}\vet{F}\cdot\dif\vet{r}.
\]
A hipótese torna $\phi(X)$ independente da escolha. Se $U$ é convexo, o segmento de $A$ a $X$ está em $U$ e pode ser a curva escolhida. No caso apenas conexo, usa-se uma poligonal, ou qualquer curva $C^1$ por pedaços, de $A$ até $X$.

Falta provar que $\partial_x\phi=P_1$, e o análogo para as outras variáveis. Seja $X=(x,y,z)\in U$. Como $U$ é aberto, existe $h\neq 0$ tal que o segmento $\sigma$ de $(x,y,z)$ até $(x+h,y,z)$ está em $U$. A concatenação de $\gamma_X$ com $\sigma$ liga $A$ a $(x+h,y,z)$. Pela independência de caminho,
\[
\phi(x+h,y,z)-\phi(x,y,z)=\int_\sigma\vet{F}\cdot\dif\vet{r}.
\]
Parametrize $\sigma(t)=(x+t,y,z)$ para $t$ entre $0$ e $h$. Então $\sigma'(t)=(1,0,0)$ e
\[
\phi(x+h,y,z)-\phi(x,y,z)
=\int_0^h P_1(x+t,y,z)\,\dif t,
\]
fórmula que vale para $h>0$ e para $h<0$. O quociente de diferença é a média
\[
\frac{\phi(x+h,y,z)-\phi(x,y,z)}{h}
=\frac{1}{h}\int_0^h P_1(x+t,y,z)\,\dif t.
\]
Seja $\eps>0$. Pela continuidade de $P_1$ em $(x,y,z)$, existe $\delta>0$ tal que $|t|<\delta$ implica $|P_1(x+t,y,z)-P_1(x,y,z)|<\eps$. Se $0<|h|<\delta$, o intervalo de extremos $0$ e $h$ tem comprimento $|h|$, e o integrando $P_1(x+t,y,z)-P_1(x,y,z)$ tem módulo menor que $\eps$. Portanto
\[
\Biggl|\int_0^h\bigl(P_1(x+t,y,z)-P_1(x,y,z)\bigr)\,\dif t\Biggr|
\le\eps\,|h|,
\]
e o quociente por $h$ dista de $P_1(x,y,z)$ menos que $\eps$. Logo $\partial_x\phi(x,y,z)=P_1(x,y,z)$.

O mesmo segmento, agora paralelo ao eixo $y$ ou ao eixo $z$, dá $\partial_y\phi=P_2$ e $\partial_z\phi=P_3$. As parciais existem numa vizinhança e coincidem com as componentes contínuas de $\vet{F}$, portanto $\phi$ é de classe $C^1$ e $\nabla\phi=\vet{F}$.
\end{proof}

A hipótese usada na construção é a continuidade de $\vet{F}$, duas vezes: para a integral de linha existir ao longo da poligonal, e para a média no segmento curto convergir ao valor no ponto. Não se usou derivada de $\vet{F}$.

\begin{observacao}
A independência de caminho equivale a dizer que $\int_C\vet{F}\cdot\dif\vet{r}=0$ para toda curva fechada $C^1$ por pedaços em $U$. Se $\gamma_1$ e $\gamma_2$ vão de $A$ a $B$, a concatenação de $\gamma_1$ com o oposto de $\gamma_2$ é fechada, e a integral dessa concatenação é a diferença das duas integrais.
\end{observacao}

\section{Conservativo implica rotacional nulo}

Esta é a direção fácil do teste. O rotacional mede a circulação infinitesimal. Um gradiente não circula, porque a integral em qualquer curva fechada pequena já é zero pelo Teorema~\ref{thm:gradiente}. A prova abaixo não estima essa circulação: ela só usa que as derivadas mistas do potencial comutam.

\base{Teorema de Clairaut: se as derivadas parciais de segunda ordem existem numa vizinhança e são contínuas, então $\partial_x\partial_y\phi=\partial_y\partial_x\phi$, e o mesmo para os outros pares. Se $\vet{F}=\nabla\phi$ e $\vet{F}$ é de classe $C^1$, as segundas parciais de $\phi$ são as primeiras parciais das componentes de $\vet{F}$, logo existem e são contínuas.}

\begin{definicao}
Seja $\vet{F}=(P,Q,R)$ de classe $C^1$ num aberto de $\R^3$. O \emph{rotacional} de $\vet{F}$ é
\[
\nabla\times\vet{F}
=\bigl(
\partial_y R-\partial_z Q,\;
\partial_z P-\partial_x R,\;
\partial_x Q-\partial_y P
\bigr).
\]
Se $\vet{F}=(P,Q)$ é de classe $C^1$ num aberto do plano, o rotacional escalar é o número $\partial_x Q-\partial_y P$, que é a terceira componente de $\nabla\times(P,Q,0)$.
\end{definicao}

\begin{teorema}
\label{thm:facil}
Seja $U$ aberto de $\R^2$ ou de $\R^3$ e seja $\vet{F}$ de classe $C^1$ em $U$. Se $\vet{F}$ é conservativo em $U$, então $\nabla\times\vet{F}=\vet{0}$ em $U$. No plano, isso é $\partial_x Q-\partial_y P=0$.
\end{teorema}

\begin{proof}
Seja $\phi$ um potencial de classe $C^1$ com $\vet{F}=\nabla\phi$. Como $\vet{F}$ é $C^1$, $\phi$ é $C^2$ e as segundas parciais são contínuas, de modo que Clairaut se aplica.

No espaço, $P=\partial_x\phi$, $Q=\partial_y\phi$ e $R=\partial_z\phi$. As três componentes do rotacional são
\begin{align*}
\partial_y R-\partial_z Q
&=\partial_y\partial_z\phi-\partial_z\partial_y\phi=0,\\
\partial_z P-\partial_x R
&=\partial_z\partial_x\phi-\partial_x\partial_z\phi=0,\\
\partial_x Q-\partial_y P
&=\partial_x\partial_y\phi-\partial_y\partial_x\phi=0.
\end{align*}
No plano, $P=\partial_x\phi$ e $Q=\partial_y\phi$, e a única conta é a terceira: $\partial_x Q-\partial_y P=\partial_x\partial_y\phi-\partial_y\partial_x\phi=0$.
\end{proof}

A hipótese de continuidade das segundas parciais de $\phi$ não é um enfeite. Ela é herdada de $\vet{F}$ ser $C^1$, e é exatamente a hipótese de Clairaut que autoriza trocar a ordem. Sem essa troca, a componente $\partial_x Q-\partial_y P$ não precisaria ser zero.

\section{O plano menos a origem}

O Teorema~\ref{thm:facil} não tem recíproca automática. O rotacional nulo em cada ponto controla a circulação em retângulos pequenos. Para passar do retângulo pequeno a um potencial no domínio inteiro, todo laço do domínio precisa poder ser contraído a um ponto sem sair do domínio. O plano sem a origem não tem essa propriedade: a circunferência unitária dá uma volta no furo e não encolhe até um ponto dentro do domínio.

\begin{contraexemplo}
\label{ex:angular}
Seja $U=\R^2\setminus\{\vet{0}\}$ e
\[
\vet{F}
=\Biggl(-\frac{y}{x^2+y^2},\;\frac{x}{x^2+y^2}\Biggr)
=(P,Q).
\]
O campo é de classe $C^\infty$ em $U$, porque o denominador não se anula fora da origem. Apesar disso, $\vet{F}$ não é conservativo em $U$.
\end{contraexemplo}

Primeiro o rotacional, por derivação direta. Escreva $P=-y(x^2+y^2)^{-1}$ e $Q=x(x^2+y^2)^{-1}$. Então
\begin{align*}
\partial_x Q
&=(x^2+y^2)^{-1}+x\cdot(-1)(x^2+y^2)^{-2}\cdot 2x
=\frac{x^2+y^2-2x^2}{(x^2+y^2)^2}
=\frac{y^2-x^2}{(x^2+y^2)^2},\\[0.4em]
\partial_y P
&=-(x^2+y^2)^{-1}+(-y)\cdot(-1)(x^2+y^2)^{-2}\cdot 2y
=\frac{-(x^2+y^2)+2y^2}{(x^2+y^2)^2}
=\frac{y^2-x^2}{(x^2+y^2)^2}.
\end{align*}
Logo $\partial_x Q-\partial_y P=0$ em todo ponto de $U$.

Agora a circulação na circunferência $x^2+y^2=1$, orientada no sentido anti-horário por $\vet{r}(t)=(\cos t,\sin t)$, $t\in[0,2\pi]$. Sobre a curva, $x^2+y^2=1$, de modo que
\[
\vet{F}(\vet{r}(t))=(-\sin t,\cos t),\qquad\vet{r}'(t)=(-\sin t,\cos t).
\]
O produto escalar vale $1$, e
\[
\oint_{x^2+y^2=1}\vet{F}\cdot\dif\vet{r}
=\int_0^{2\pi}1\,\dif t=2\pi.
\]
Se existisse um potencial em $U$, o Teorema~\ref{thm:gradiente} daria circulação zero. Como $2\pi\neq 0$, $\vet{F}$ não é conservativo em $\R^2\setminus\{\vet{0}\}$.

\cuidado{A recíproca de ``conservativo implica rotacional nulo'' pede domínio simplesmente conexo: todo laço se contrai a um ponto dentro do domínio. No plano menos a origem, o campo acima tem $\partial_x Q-\partial_y P=0$ em cada ponto e mesmo assim $\oint_{x^2+y^2=1}\vet{F}\cdot\dif\vet{r}=2\pi\neq 0$. O furo guarda uma volta que o rotacional, calculado fora da origem, não vê. Em $\R^2$ menos um raio --- uma semi-reta fechada que parte da origem --- o domínio restante é simplesmente conexo, e a mesma fórmula é o gradiente do argumento $\theta$.}

\begin{proposicao}
\label{prop:theta}
Seja $U=\R^2\setminus\{(x,0):x\le 0\}$, o plano menos a semi-reta dos $x$ não positivos. Seja $\theta\colon U\to(-\pi,\pi)$ o argumento, isto é, o ângulo do eixo $x$ positivo até o ponto $(x,y)$. Então $\theta$ é de classe $C^\infty$ em $U$ e $\nabla\theta=\vet{F}$ em $U$, com $\vet{F}$ o campo do Contraexemplo~\ref{ex:angular}.
\end{proposicao}

\begin{proof}
A semi-reta removida é exatamente o corte onde o argumento saltaria de $\pi$ para $-\pi$. Em $U$, $\theta$ fica em $(-\pi,\pi)$ e não salta. Três fórmulas cobrem $U$, cada uma $C^\infty$ no seu aberto:
\begin{align*}
x>0&\colon & \theta&=\arctan(y/x),\\
y>0&\colon & \theta&=\pi/2-\arctan(x/y),\\
y<0&\colon & \theta&=-\pi/2-\arctan(x/y).
\end{align*}
Todo ponto de $U$ está em pelo menos um desses abertos, e as fórmulas coincidem nas interseções porque todas descrevem o mesmo ângulo em $(-\pi,\pi)$.

No semiplano $x>0$,
\begin{align*}
\partial_x\theta
&=\frac{1}{1+(y/x)^2}\cdot\Bigl(-\frac{y}{x^2}\Bigr)
=\frac{x^2}{x^2+y^2}\cdot\Bigl(-\frac{y}{x^2}\Bigr)
=-\frac{y}{x^2+y^2},\\[0.3em]
\partial_y\theta
&=\frac{1}{1+(y/x)^2}\cdot\frac{1}{x}
=\frac{x^2}{x^2+y^2}\cdot\frac{1}{x}
=\frac{x}{x^2+y^2}.
\end{align*}
No semiplano $y>0$,
\begin{align*}
\partial_x\theta
&=-\frac{1}{1+(x/y)^2}\cdot\frac{1}{y}
=-\frac{y^2}{x^2+y^2}\cdot\frac{1}{y}
=-\frac{y}{x^2+y^2},\\[0.3em]
\partial_y\theta
&=-\frac{1}{1+(x/y)^2}\cdot\Bigl(-\frac{x}{y^2}\Bigr)
=\frac{y^2}{x^2+y^2}\cdot\frac{x}{y^2}
=\frac{x}{x^2+y^2}.
\end{align*}
No semiplano $y<0$ a fórmula difere da anterior por uma constante, $-\pi$ no lugar de $\pi/2$ diante do mesmo $\arctan(x/y)$, então as parciais são as mesmas. Portanto $\nabla\theta=\vet{F}$ em $U$.
\end{proof}

O domínio $U$ é simplesmente conexo na linguagem do curso: um laço em $U$ não consegue dar a volta na origem, porque a semi-reta removida barra esse giro, e o laço se contrai a um ponto sem sair de $U$. A circunferência unitária inteira não cabe em $U$. A circulação $2\pi$ vive no plano furado; no plano cortado, o Teorema~\ref{thm:gradiente} devolve circulação zero para toda curva fechada.

\section{O teste no retângulo}

O teste prático junta as duas direções. A ida é o Teorema~\ref{thm:facil} e vale em qualquer aberto. A volta pede topologia.

\begin{teorema}[Teste prático]
\label{thm:teste}
Seja $U\subset\R^2$ aberto, e suponha que toda curva fechada em $U$ se contrai a um ponto por uma deformação contínua que permanece em $U$. Se $\vet{F}=(P,Q)$ é de classe $C^1$ em $U$ e $\partial_x Q-\partial_y P=0$ em $U$, então $\vet{F}$ é conservativo em $U$.

Seja $B\subset\R^3$ uma bola aberta. Se $\vet{F}$ é de classe $C^1$ em $B$ e $\nabla\times\vet{F}=\vet{0}$ em $B$, então $\vet{F}$ é conservativo em $B$.
\end{teorema}

Uma bola, um retângulo aberto, o plano inteiro e o plano menos um raio entram na hipótese: todo laço se contrai a um ponto dentro do domínio. O plano menos a origem não entra. A prova geral, reduzindo um laço arbitrário a laços pequenos por essa contração, não cabe nesta folha. O que cabe é o retângulo, e a conta é direta: constrói-se o potencial por um caminho em ``L'' e usa-se $\partial_x Q=\partial_y P$ para identificar as parciais. Não se invoca o teorema de Green.

\begin{teorema}
\label{thm:retangulo}
Seja $R=(\alpha,\beta)\times(\gamma,\delta)$ um retângulo aberto e seja $\vet{F}=(P,Q)$ de classe $C^1$ em $R$, com $\partial_x Q=\partial_y P$ em $R$. Então $\vet{F}$ é conservativo em $R$.
\end{teorema}

\begin{proof}
Fixe $(x_0,y_0)\in R$. Para $(x,y)\in R$, os dois segmentos de $(x_0,y_0)$ até $(x,y_0)$ e de $(x,y_0)$ até $(x,y)$ ficam em $R$, porque $R$ é um produto. Defina
\[
\phi(x,y)
=\int_{x_0}^{x} P(t,y_0)\,\dif t
+\int_{y_0}^{y} Q(x,s)\,\dif s.
\]
A primeira integral não depende de $y$. Pelo teorema fundamental do cálculo, $\partial_y\phi(x,y)=Q(x,y)$.

Para a outra parcial, o primeiro termo contribui $P(x,y_0)$. No segundo, os limites não dependem de $x$. Afirmamos que
\[
\partial_x\int_{y_0}^{y} Q(x,s)\,\dif s
=\int_{y_0}^{y}\partial_x Q(x,s)\,\dif s.
\]
De fato, seja $\psi(x)=\int_{y_0}^{y} Q(x,s)\,\dif s$ e seja $I$ o intervalo de extremos $y_0$ e $y$. Pelo teorema do valor médio, para cada $s\in I$ existe $\xi$ entre $x$ e $x+h$ com
\[
\frac{Q(x+h,s)-Q(x,s)}{h}=\partial_x Q(\xi,s).
\]
Como $\partial_x Q$ é contínua num compacto em volta do segmento $\{x\}\times I$, ela é uniformemente contínua nesse compacto. Dado $\eps>0$, para $|h|$ pequeno vale $|\partial_x Q(\xi,s)-\partial_x Q(x,s)|<\eps$ para todo $s\in I$, qualquer que seja o $\xi$ intermediário. Então
\[
\Biggl|\frac{\psi(x+h)-\psi(x)}{h}-\int_{y_0}^{y}\partial_x Q(x,s)\,\dif s\Biggr|
\le\eps\,|y-y_0|.
\]
Isso prova a derivação sob o sinal de integral. Portanto
\begin{align*}
\partial_x\phi(x,y)
&=P(x,y_0)+\int_{y_0}^{y}\partial_x Q(x,s)\,\dif s
=P(x,y_0)+\int_{y_0}^{y}\partial_y P(x,s)\,\dif s\\
&=P(x,y_0)+\bigl(P(x,y)-P(x,y_0)\bigr)
=P(x,y),
\end{align*}
onde a penúltima igualdade é o teorema fundamental do cálculo na variável $s$, aplicado a $s\mapsto P(x,s)$, cuja derivada $\partial_y P(x,s)$ é contínua. As parciais de $\phi$ são $P$ e $Q$, contínuas, logo $\phi$ é $C^1$ e $\nabla\phi=\vet{F}$.
\end{proof}

A mesma construção, com uma aresta horizontal e uma vertical, é o que o plano furado não permite de modo global: o ``L'' de um ponto a outro pode ser obrigado a pular a origem, e a igualdade $\partial_x Q=\partial_y P$ fora da origem não cola esses pedaços num potencial único.

\section{Exercícios}

\begin{exercicio}[{[Básico --- campos e unidades]}]
Classifique cada campo e dê a unidade.
\begin{enumerate}
\item $T(x,y,z)$, temperatura em kelvin.
\item $\vet{v}(x,y,z)$, velocidade de um fluido em metro por segundo.
\item $-\nabla T$, se $T$ é de classe $C^1$ e está em kelvin, com as coordenadas em metros.
\end{enumerate}
Diga, em cada caso, qual das duas integrais de linha faz sentido físico imediato: $\int f\,\dif s$ ou $\int\vet{F}\cdot\dif\vet{r}$.
\end{exercicio}

\begin{exercicio}[{[Básico --- massa de um fio]}]
Um fio percorre $\vet{r}(t)=(t,t^2)$, $t\in[0,1]$, com densidade linear $\delta(x,y)=2x$. Calcule a massa $\int_C\delta\,\dif s$.
\end{exercicio}

\begin{exercicio}[{[Básico --- trabalho e orientação]}]
Seja $\vet{F}=(y,x)$. Calcule $\int_C\vet{F}\cdot\dif\vet{r}$ no segmento de $(0,0)$ a $(1,1)$, parametrizado por $\vet{r}(t)=(t,t)$, $t\in[0,1]$. Calcule a integral no sentido oposto. Exiba um potencial e confira os dois valores por $\phi(B)-\phi(A)$.
\end{exercicio}

\begin{exercicio}[{[Básico --- integral escalar e o sentido do percurso]}]
Seja $C$ a circunferência unitária. Calcule $\int_C 1\,\dif s$ no sentido anti-horário e no sentido horário. Calcule também $\int_C x\,\dif s$ nos dois sentidos. Escreva a parametrização do sentido horário e mostre, nessa conta, o passo em que a norma absorve o sinal.
\end{exercicio}

\begin{exercicio}[{[Intermediário --- potencial e trabalho]}]
Seja $\vet{F}=(2xy+1,\, x^2+\cos y)$ em $\R^2$.
\begin{enumerate}
\item Verifique $\partial_x Q=\partial_y P$.
\item Determine um potencial.
\item Calcule o trabalho de $(0,0)$ até $(1,\pi/2)$.
\end{enumerate}
Qual hipótese do teste prático está sendo usada ao concluir que o potencial serve em todo o plano?
\end{exercicio}

\begin{exercicio}[{[Intermediário --- o campo de rotação]}]
Seja $\vet{F}=(-y,x)$ em $\R^2$.
\begin{enumerate}
\item Calcule $\partial_x Q-\partial_y P$.
\item Calcule a circulação na circunferência unitária anti-horária.
\item Conclua, de duas maneiras, que $\vet{F}$ não é conservativo em $\R^2$.
\item A forma angular $\bigl(-y/(x^2+y^2),\, x/(x^2+y^2)\bigr)$ também não é conservativa no plano furado. O obstáculo é o mesmo deste item, ou é outro?
\end{enumerate}
\end{exercicio}

\begin{exercicio}[{[Intermediário --- potencial pela receita do retângulo]}]
Seja $\vet{F}=(e^x\sin y,\, e^x\cos y)$. Use a construção do Teorema~\ref{thm:retangulo} a partir de $(0,0)$:
\[
\phi(x,y)=\int_0^x P(t,0)\,\dif t+\int_0^y Q(x,s)\,\dif s.
\]
Calcule $\phi$ e verifique $\nabla\phi=\vet{F}$.
\end{exercicio}

\begin{exercicio}[{[Avançado --- independência de caminho]}]
Seja $U\subset\R^n$ aberto e conexo e seja $\vet{F}\colon U\to\R^n$ contínuo. Suponha que $\int_C\vet{F}\cdot\dif\vet{r}=0$ para toda curva fechada $C^1$ por pedaços em $U$.
\begin{enumerate}
\item Prove que a integral entre dois pontos não depende do caminho.
\item Fixe $A$ e defina $\phi(X)$ pela integral de $A$ até $X$. Prove que $\partial_x\phi=P_1$, escrevendo o acréscimo ao longo de um segmento horizontal pequeno e estimando a média pela continuidade de $P_1$.
\end{enumerate}
\end{exercicio}

\begin{exercicio}[{[Avançado --- o argumento no plano cortado]}]
Seja $U=\R^2\setminus\{(x,0):x\le 0\}$ e seja $\theta\in(-\pi,\pi)$ o argumento.
\begin{enumerate}
\item No semiplano $x>0$, com $\theta=\arctan(y/x)$, calcule $\partial_x\theta$ e $\partial_y\theta$.
\item No semiplano $y>0$, com $\theta=\pi/2-\arctan(x/y)$, calcule as mesmas parciais.
\item Conclua que a forma angular é conservativa em $U$ e não é conservativa em $\R^2\setminus\{\vet{0}\}$. Aponte, no segundo domínio, o laço que não se contrai a um ponto.
\end{enumerate}
\end{exercicio}

\end{document}
